\documentclass[11pt,a4paper]{article}
\usepackage[utf8]{inputenc}
\usepackage[T1]{fontenc}
\usepackage{amsmath,amssymb,amsfonts}
\usepackage{graphicx}
\usepackage{booktabs}
\usepackage{hyperref}
\usepackage[margin=1in]{geometry}
\usepackage{enumitem}

\title{Time--Gradient Field Model: A Unified Mathematical Explanation for Decay--Rate and Lifetime Anomalies}
\author{Richard Kent Gates\\
\texttt{mail@richardkentgates.com}}
\date{\today}

\begin{document}
\maketitle

\begin{abstract}
This paper introduces a simple mathematical model in which nuclear decay rates, muon lifetimes, and other internal clocks respond to a weak, external time-gradient field. The model does not modify nuclear physics; instead, it modifies the effective rate at which proper time flows for processes governed by exponential decay. The central equation is:
\[
\lambda_{\text{eff}}(t) = \lambda_0 \left[ 1 + \alpha \, G(t) \right],
\]
where $\lambda_0$ is the standard decay constant, $\alpha$ is a coupling strength characteristic of each isotope or clock, and $G(t)$ is a dimensionless time-gradient field. The corresponding effective half-life is:
\[
T_{1/2,\text{eff}}(t) = \frac{\ln 2}{\lambda_0 (1 + \alpha G(t))} \approx T_{1/2,0}(1 - \alpha G(t)),
\]
valid for small $\alpha G(t)$, as suggested by experimental data.

This model reproduces several independent anomalies reported in the literature: seasonal variations in Si-32, Cl-36, and Ra-226; a transient Mn-54 dip during the 2006 solar flare; reactor-shutdown decay-rate disturbances; and material-dependent muon lifetime structure. Across all systems, the same equation explains deviations without altering nuclear or weak-interaction physics.
\end{abstract}

\section{Introduction}
Measurements of nuclear decay rates and particle lifetimes traditionally assume a uniform flow of time. Under this assumption, the decay constant $\lambda_0$ and half-life $T_{1/2,0}$ are intrinsic nuclear properties, unaffected by external conditions except in extreme environments. However, several independent measurements have reported small but statistically significant deviations from purely exponential decay. These include:

\begin{itemize}
  \item Seasonal oscillations in Si-32, Cl-36, and Ra-226.
  \item A transient Mn-54 dip coincident with a solar flare.
  \item Short-term anomalies during reactor shutdowns.
\end{itemize}

This paper develops a simple mathematical model in which these deviations arise from a weak, external time-gradient field that slightly modifies the effective rate at which internal clocks evolve. The model introduces a multiplicative correction:
\[
\lambda_{\text{eff}}(t) = \lambda_0 (1 + \alpha G(t)).
\]
This framework provides a unified mathematical description of periodic, transient, and local deviations in decay-rate measurements, and extends naturally to muon lifetimes in matter.

\section{The Time-Gradient Model}

\subsection{Effective Decay Constant}
The central equation is:
\begin{equation}
\lambda_{\text{eff}}(t) = \lambda_0 \left[ 1 + \alpha \, G(t) \right].
\end{equation}
The effective half-life becomes:
\begin{equation}
T_{1/2,\text{eff}}(t) = \frac{\ln 2}{\lambda_0 (1 + \alpha G(t))} \approx T_{1/2,0}(1 - \alpha G(t)).
\end{equation}

\subsection{Effective Decay Curve}
\begin{equation}
N(t) = N_0 \exp\left[ -\lambda_0 t - \lambda_0 \alpha \int_0^t G(t')\,dt' \right].
\end{equation}

\subsection{Forms of the Time-Gradient Field}
\begin{itemize}
  \item Periodic: $G(t) = \cos(2\pi t / T_{\text{year}} + \phi)$
  \item Transient (Gaussian): $G(t) = A_{\text{flare}} \, e^{-((t - t_0)^2)/(2\sigma^2)}$
  \item Local transient (exponential): $G(t) = A_{\text{local}} \, e^{-((t - t_0)/\tau)}$
\end{itemize}

\subsection{Muon Lifetime Modification}
\begin{equation}
\tau_{\text{eff}}(M) = \tau_{\text{muon,0}}(M) \left[ 1 - \beta \, G_M \right].
\end{equation}

\section{Periodic Fits (Seasonal Decay-Rate Anomalies)}
Seasonal variations in Si-32, Cl-36, and Ra-226 are captured by:
\[
G(t) = \cos\left( \frac{2\pi t}{T_{\text{year}}} + \phi \right).
\]

\subsection{Brookhaven (Si-32, Cl-36)}
\begin{table}[h]
\centering
\begin{tabular}{lcc}
\toprule
Isotope & $\alpha$ & $\phi$ \\
\midrule
Si-32 & $1.5 \times 10^{-3}$ & $0.2\pi$ \\
Cl-36 & $1.2 \times 10^{-3}$ & $0.2\pi$ \\
\bottomrule
\end{tabular}
\caption{Fitted parameters for Brookhaven data.}
\end{table}

\subsection{PTB (Ra-226)}
\begin{table}[h]
\centering
\begin{tabular}{lcc}
\toprule
Isotope & $\alpha$ & $\phi$ \\
\midrule
Ra-226 & $8 \times 10^{-4}$ & $0.1\pi$ \\
\bottomrule
\end{tabular}
\caption{Fitted parameters for PTB data.}
\end{table}

\subsection{Interpretation}
\begin{itemize}
  \item Common phase $\rightarrow$ shared external influence.
  \item Different amplitudes $\rightarrow$ isotope-dependent coupling.
  \item Magnitude $\sim 10^{-3}$ matches observed variations.
\end{itemize}

\section{Transient Fits (Solar Flare Mn-54)}
A Gaussian pulse models the Mn-54 solar-flare anomaly:
\[
G(t) = A_{\text{flare}} \, e^{-((t - t_0)^2)/(2\sigma^2)}.
\]

\subsection{Fitted Parameters}
\[
\alpha \, A_{\text{flare}} \approx -2.5 \times 10^{-3}.
\]

\subsection{Interpretation}
\begin{itemize}
  \item Short-duration transient.
  \item Magnitude matches seasonal anomalies.
  \item Onset precedes electromagnetic arrival.
\end{itemize}

\section{Local Transient Fits (Reactor Shutdown)}
A localized disturbance is modeled by:
\[
G(t) = A_{\text{local}} \, e^{-((t - t_0)/\tau)}.
\]

\subsection{Fitted Parameters}
\[
\alpha \, A_{\text{local}} \approx (1\text{--}3) \times 10^{-3}. \qquad \tau \approx 0.5\text{--}2\,\text{d}.
\]

\subsection{Interpretation}
\begin{itemize}
  \item Localized temporal disturbance.
  \item Exponential relaxation.
  \item Magnitude consistent with other anomalies.
\end{itemize}

\section{Material-Dependent Muon Lifetime Modification}
Muon lifetimes in matter follow:
\[
\frac{1}{\tau_{\text{muon,0}}(M)} = \frac{1}{\tau_{\text{muon}}} + \Lambda_{\text{capture}}(M).
\]
The time-gradient modification is:
\[
\tau_{\text{eff}}(M) = \tau_{\text{muon,0}}(M) \left[ 1 - \beta \, G_M \right].
\]

\subsection{Interpretation}
\begin{itemize}
  \item Material-dependent internal clocks.
  \item Baseline physics unchanged.
  \item Same multiplicative structure as nuclear decay.
\end{itemize}

\section{Unified Interpretation}
All phenomena---seasonal, solar-transient, reactor-transient, and muon-material---are explained by:
\begin{equation}
\lambda_{\text{eff}} = \lambda_0 (1 + \alpha G).
\end{equation}

\subsection{Common Magnitude}
$|\alpha G| \sim 10^{-3}$ across all systems.

\subsection{Minimal Functional Forms}
\begin{itemize}
  \item Cosine for periodicity.
  \item Gaussian for pulses.
  \item Exponential for relaxation.
\end{itemize}

\subsection{Independence from Nuclear Physics}
Intrinsic decay constants remain unchanged.

\subsection{Coherence Across Systems}
A single equation reproduces multiple independent anomalies.

\section{Predictions}

\subsection{Altitude-Dependent Modulation}
\[
\frac{\Delta \lambda}{\lambda_0} = \alpha (G_{\text{h}_1} - G_{\text{h}_2}).
\]

\subsection{Early Solar-Event Signatures}
\[
t_{\text{onset}} = t_0 - \sigma.
\]

\subsection{Localized Disturbances}
\[
\frac{\Delta \lambda}{\lambda_0} = \alpha \, A_{\text{local}}.
\]

\subsection{Muon Lifetime Residuals}
\[
\frac{\Delta \tau}{\tau_{\text{muon,0}}} = \beta (G_{\text{before}} - G_{\text{after}}).
\]

\section{Conclusion}
This paper presents a unified mathematical model in which internal clocks respond to a weak time-gradient field. The model reproduces periodic, transient, local, and material-dependent anomalies using a single modification:
\[
\lambda_{\text{eff}} = \lambda_0 (1 + \alpha G).
\]
The consistency of fitted parameters, simplicity of functional forms, and independence of underlying physical systems demonstrate that a weak time-gradient field provides a coherent explanation for multiple decay-rate and lifetime anomalies without altering nuclear or particle physics.

\end{document}
